\section*{Chapter \ref{ch:g10:production-possibility-frontier} --- The Production Possibility Frontier}
% figdata: figdata/grade-10/02-production-possibility-frontier/out/ (war_usa, war_gbr); Brenland's
% valleys and regions are invented, with their numbers in the chapter.

\begin{solution}{exo:g10:production-possibility-frontier:1}
United States, 1944: output per person 180.2, consumption per person 110.0 (1939 = 100). United
Kingdom, 1943: output per person 123.7, consumption per person 83.5.
\end{solution}

\begin{solution}{exo:g10:production-possibility-frontier:2}
In 1946 (128.4). Output per person fell from 180.2 in 1944 to 150.4 in 1947, a fall of $1 - 150.4/180.2
\approx 16.5\,\%$, as war production stopped; consumption per person rose at the same time.
\end{solution}

\begin{solution}{exo:g10:production-possibility-frontier:3}
(a) $3 \times 100 + 2 \times 450 = 1200$: on the frontier. (b) $900 + 400 = 1300 > 1200$: out of reach.
(c) $150 + 1000 = 1150 < 1200$: inside (50 hours unused).
\end{solution}

\begin{solution}{exo:g10:production-possibility-frontier:4}
With $F$ baskets of fish and $R$ of fruit: $2F + 3R = 600$, so $R = 200 - \tfrac{2}{3}F$. At most
$600/2 = 300$ baskets of fish, or $600/3 = 200$ of fruit.
\end{solution}

\begin{solution}{exo:g10:production-possibility-frontier:5}
A basket of fish takes 2 hours, which would make $2/3 \approx 0.67$ basket of fruit: that is its
opportunity cost. A basket of fruit costs $3/2 = 1.5$ baskets of fish.
\end{solution}

\begin{solution}{exo:g10:production-possibility-frontier:6}
East: at most 300 fish or 200 fruit; a basket of fish costs $2/3$ basket of fruit. West: at most 150
fish or 300 fruit; a basket of fish costs $4/2 = 2$ baskets of fruit. Together: at most $300 + 150 =
450$ baskets of fish, or $200 + 300 = 500$ of fruit. From all fruit, the East switches to fish first (it
gives up less fruit per basket of fish): 300 fish for 200 fruit, reaching the kink $(300, 300)$ (fish,
fruit). Then the West: 150 more fish for 300 fruit, reaching $(450, 0)$. The frontier joins $(0, 500)$,
$(300, 300)$ and $(450, 0)$, with slopes $-2/3$ then $-2$: bowed outward.
\end{solution}

\begin{solution}{exo:g10:production-possibility-frontier:7}
From 100 to 300 tonnes the grain comes from the North Valley, at 1.5 rolls a tonne: $200 \times 1.5 =
300$ rolls. From 400 to 500 tonnes it comes from the South Valley, at 6 rolls a tonne: $100 \times 6 =
600$ rolls. Half as much grain costs twice as much cloth, because beyond 400 tonnes the grain must be
grown where it is costly.
\end{solution}

\begin{solution}{exo:g10:production-possibility-frontier:8}
The South Valley can now make $1200/4 = 300$ tonnes of grain instead of 200, so the grain end moves from
600 to 700 tonnes. The cloth end does not move (1800 rolls). In the South a tonne of grain now costs
$4/1 = 4$ rolls, still more than the North's 1.5, so the North still switches first and the kink stays at
$(400, 1200)$. Only the second piece changes: it now runs from $(400, 1200)$ to $(700, 0)$, with slope
$-4$ instead of $-6$.
\end{solution}

\begin{solution}{exo:g10:production-possibility-frontier:9}
North: 200 tonnes take 600 hours; the other 600 make $600/2 = 300$ rolls. South: 100 tonnes take 600
hours; the other 600 make 600 rolls. Total: 900 rolls. On the frontier, 300 tonnes all from the North
leave $1800 - 1.5 \times 300 = 1350$ rolls. The first combination is inside the frontier: not \omterm{def:g10:production-possibility-frontier:productive-efficiency}{productively
efficient}; 450 rolls are lost by growing grain in the South.
\end{solution}

\begin{solution}{exo:g10:production-possibility-frontier:10}
From $(10, 95)$ to $(20, 85)$: 10 units of cloth for 10 of grain, a cost of 1 per unit. From $(20, 85)$
to $(30, 70)$: 15 for 10, a cost of 1.5 per unit. The cost rises as more grain is produced, as the \omterm{def:g10:production-possibility-frontier:increasing-cost}{law of
increasing opportunity cost} says.
\end{solution}

\begin{solution}{exo:g10:production-possibility-frontier:11}
United States: from 17.2\,\% to 1.2\,\%, a fall of 16.0 points. United Kingdom: from 5.8\,\% to 0.4\,\%
(1943 and 1944), a fall of 5.4 points. The United States had many idle workers to put to work, so it
could raise war output without taking workers from consumer goods: it started inside its frontier. The
United Kingdom had few, so its war output came mostly from moving workers and factories out of consumer
goods: a move along the frontier.
\end{solution}

\begin{solution}{exo:g10:production-possibility-frontier:12}
An observation that each extra unit of food cost the same amount of other goods as the one before, or
less: for example, three years in which food output rose by equal steps and output of the other goods
fell by equal steps (a straight frontier), or by smaller and smaller steps (a frontier bowed inward). The
factors and technology must have stayed the same because the claim is about the shape of one frontier: if
the frontier moved outward at the same time, the other goods could fall by less for reasons that have
nothing to do with the cost of food.
\end{solution}

\begin{solution}{exo:g10:production-possibility-frontier:13}
Let both regions give up $k$ rolls of cloth for each tonne of grain, and let $C_{\max}$ be the most cloth
they make together. Whichever region grows it, each tonne of grain costs $k$ rolls, so producing $G$ tonnes
in any way that keeps all hours in use leaves $C = C_{\max} - kG$, up to the most grain the two can make.
This is a straight line of slope $-k$: there is no point where the cost changes, so no kink.
\end{solution}

\begin{solution}{exo:g10:production-possibility-frontier:14}
West Valley: at most $600/2 = 300$ tonnes of grain or $600/3 = 200$ rolls; a tonne of grain costs $2/3$
roll. The three valleys together: at most $600 + 1200 + 200 = 2000$ rolls (no grain) or $400 + 200 + 300
= 900$ tonnes (no cloth). Order of switching to grain: West ($2/3$), North (1.5), South (6). From
$(0, 2000)$: the West's 300 tonnes cost 200 rolls, giving the kink $(300, 1800)$; the North's 400 tonnes
cost 600 rolls, giving the kink $(700, 1200)$; the South's 200 tonnes cost \num{1200} rolls, ending at
$(900, 0)$. Slopes: $-2/3$, $-1.5$, $-6$.
\end{solution}

\begin{solution}{exo:g10:production-possibility-frontier:15}
(a) This year $c = 100$; next year the frontier is unchanged, so $c = 100$: 200 in all. (b) This year $c =
100 - 2 \times 20 = 60$; next year the frontier is $c + 2m = 100 + 3 \times 20 = 160$, so with no machines
$c = 160$: 220 in all. Plan (b) gives 20 more consumer goods over the two years, at the price of 40 fewer
this year.
\end{solution}

\begin{solution}{pb:g10:production-possibility-frontier:1}
Quantities are in thousands of tonnes (grain) and thousands of metres (cloth).

\textbf{1.} Hills: $300/10 = 30$ of grain or $300/2 = 150$ of cloth. River: $400/4 = 100$ or $400/4 = 100$.
Plain: $300/2 = 150$ or $300/6 = 50$.

\textbf{2.} Hills: $10/2 = 5$ metres of cloth per tonne of grain. River: $4/4 = 1$. Plain: $2/6 \approx
0.33$.

\textbf{3.} The Plain has the lowest opportunity cost of grain. The Hills have the lowest opportunity cost
of cloth: $2/10 = 0.2$ tonne of grain per metre (River 1, Plain 3).

\textbf{4.} $4G + 4C = 400$, so $C = 100 - G$.

\textbf{5.} All cloth: $150 + 100 + 50 = 300$, so $(0, 300)$. All grain: $30 + 100 + 150 = 280$, so
$(280, 0)$.

\textbf{6.} Plain, then River, then Hills: each extra tonne of grain must come from the region where it
costs least cloth (\cref{prop:g10:production-possibility-frontier:joint}).

\textbf{7.} The Plain's 150 cost $150/3 = 50$ of cloth: first kink $(150, 250)$. The River's 100 cost 100:
second kink $(250, 150)$. The Hills' 30 cost 150, ending at $(280, 0)$.

\textbf{8.} $-1/3$ from $(0, 300)$ to $(150, 250)$; $-1$ from $(150, 250)$ to $(250, 150)$; $-5$ from
$(250, 150)$ to $(280, 0)$.

\textbf{9.} Three straight pieces joining $(0, 300)$, $(150, 250)$, $(250, 150)$ and $(280, 0)$, each
steeper than the one before: a frontier bowed outward.

\textbf{10.} 200 lies on the River piece: $250 - 1 \times (200 - 150) = 200$ thousand metres of cloth. The
River region produces both goods (50 thousand tonnes of grain and 50 thousand metres of cloth); the Plain
makes only grain, the Hills only cloth.

\textbf{11.} 260 lies on the Hills piece: $150 - 5 \times (260 - 250) = 100$ thousand metres.

\textbf{12.} $200 - 100 = 100$ thousand metres of cloth, that is \num{100000} metres, for \num{60000}
tonnes of grain: on average $100/60 \approx 1.67$ metres per tonne.

\textbf{13.} The Hills make 30 of grain and no cloth; the River makes 100 of grain and no cloth; the Plain
makes $260 - 130 = 130$ of grain, using $130 \times 2 = 260$ of its 300 thousand hours, and with the other
40 makes $40/6 \approx 6.67$ thousand metres of cloth. Brenland makes about \num{6667} metres of cloth,
against \num{100000} in the efficient plan: about \num{93333} metres are wasted.

\textbf{14.} Far inside the frontier: the same grain with much less cloth. It is not \omterm{def:g10:production-possibility-frontier:productive-efficiency}{productively
efficient}, because grain is grown in the Hills, where it costs the most cloth, while the Plain still makes
cloth.

\textbf{15.} The first tonnes come from the Plain: $1000 \times \tfrac13 \approx 333$ metres of cloth.

\textbf{16.} The last tonnes come from the Hills: $1000 \times 5 = 5000$ metres of cloth.

\textbf{17.} $5 / \tfrac13 = 15$: the last tonne costs fifteen times as much cloth as the first. This is the
\omterm{def:g10:production-possibility-frontier:increasing-cost}{law of increasing opportunity cost}.

\textbf{18.} The Hills can now make $300/1 = 300$ thousand metres of cloth instead of 150, so the cloth end
moves from $(0, 300)$ to $(0, 450)$. The grain end stays at $(280, 0)$.

\textbf{19.} In the Hills a tonne of grain now costs $10/1 = 10$ metres of cloth (the Hills are still the
dearest place for grain), so the last \num{1000} tonnes cost \num{10000} metres.

\textbf{20.} Before the new loom, the last \num{1000} tonnes of grain cost \num{5000} metres of cloth:
three times the campaign's average of about \num{1667} metres per \num{1000} tonnes ($5000/1667 \approx 3.0$),
and fifteen times the first \num{1000} tonnes.
\end{solution}
